\section{Introdução}

A capacidade de extrair padrões complexos de grandes volumes de dados tornou-se um pilar fundamental para a tomada de decisões em diversas esferas da sociedade moderna. No contexto socioeconômico, a aplicação de algoritmos de Aprendizado de Máquina (\textit{Machine Learning}) permite identificar fatores determinantes que influenciam a mobilidade social e a distribuição de renda. Tais análises são cruciais não apenas para a compreensão estrutural de uma população, mas também para auxiliar na formulação de políticas públicas mais eficientes.

Dentre as diversas tarefas de aprendizado supervisionado, a classificação destaca-se pela sua capacidade de categorizar instâncias em classes predefinidas com base em um conjunto de atributos. Nesse sentido, o presente trabalho aborda o problema de classificação binária utilizando o renomado \textit{Adult Income Dataset} (também conhecido como \textit{Census Income}), extraído da base de dados do censo norte-americano de 1994 \cite{becker1996adult}. O desafio inerente a este problema consiste em prever se a renda anual de um indivíduo ultrapassa o limiar de 50 mil dólares, baseando-se em características demográficas e empregatícias, como idade, nível de escolaridade, ocupação, gênero e horas trabalhadas por semana.

Este problema apresenta desafios característicos de bases de dados reais que exigem rigor metodológico, notadamente a presença de valores ausentes e a heterogeneidade dos dados (uma mistura de atributos numéricos contínuos e nominais categóricos). O principal obstáculo, contudo, é o forte desbalanceamento entre as classes, visto que a esmagadora maioria dos registros pertence à classe de menor renda. Por essa razão, tais características demandam um pré-processamento minucioso e o uso de métricas de avaliação que penalizem adequadamente o viés do modelo em direção à classe majoritária.

Diante deste cenário, o objetivo principal deste trabalho é investigar e comparar o desempenho de diferentes algoritmos de classificação na tarefa de prever a renda dos indivíduos. Para tanto, foram avaliados seis modelos distintos: Árvore de Decisão, \textit{K-Nearest Neighbors} (KNN), \textit{Naive Bayes}, Regressão Logística, \textit{Support Vector Machine} (LinearSVC) e Redes Neurais Artificiais (MLP). A comparação é fundamentada em um rigoroso procedimento experimental que adota a otimização de hiperparâmetros e a técnica de validação cruzada estratificada (\textit{10-fold cross-validation}), garantindo assim a robustez e a confiabilidade estatística dos resultados obtidos.

Espera-se, com este estudo, evidenciar quais arquiteturas algorítmicas lidam melhor com o desbalanceamento intrínseco aos dados demográficos. Além disso, as análises aqui apresentadas buscam fornecer diretrizes práticas sobre a importância da etapa de pré-processamento e da calibragem de hiperparâmetros para a construção de modelos preditivos mais precisos e confiáveis em cenários reais.